%\section{Puschkin -- Winterabend}
%\stepcounter{ctrpoeminyear}
%%\addcontentsline{toc}{section}{\thechapter.\arabic{ctrpoeminyear} \;}

\section*{Kakinomoto -- Manyoshu-2-131}
%\clearpage
%\stepcounter{ctrpoeminyear}
\addcontentsline{toc}{section}{2026.Ü2 \;Kakinomoto -- Manyoshu-2-131}
\fancyhead[LO,RE]{2026.Ü2}

\begin{strophe}
\vers{In Iwami, in Tsuno am Meere,}
\vers{wo die Menschen mir flüsternd erklären,}
\vers{dass es dort keine Buchten gebe}
\vers{und auch keine Lagune schwebe.}
\vers{Dort, wie auch immer es zu sein vermag,}
\vers{ob es nun jemals dort Buchten gab}
\vers{oder ob dort je eine Lagune lag,}
\vers{dort jagen sie Waale den ganzen Tag.}
\end{strophe}

\begin{strophe}
\vers{Ich gehe dann gerne den Strand auf und ab,}
\vers{blicke beim Wada-Hafen aufs Wasser herab,}
\vers{sehe die Überreste des Filetierens dort schwimmen}
\vers{auf denen die wertvollen Gräser drob wimmeln.}
\vers{Blauer und küstenferner Seetang ist dort.}
\vers{Ein wertvolles Schimmern spiegelt sich fort.}
\end{strophe}

\begin{strophe}
\vers{-- Das Rasseln der Wellen am Morgen,}
\vers{durch den schwimmenden Seetang erworben --}
\vers{-- Das Rauschen der Wellen am Abend,}
\vers{wenn die Winde auf dem Wasser traben --}
\end{strophe}

\begin{strophe}
\vers{Die Wellen schwenken vor und zurück,}
\vers{verschlingen das Land und zerfallen beim Blick.}
\vers{So schwenkt der wertvolle Seetang auch mit,}
\vers{wie wir damals, als ich an deiner Seite einschlief.}
\end{strophe}

\begin{strophe}
\vers{Schwester, was nur spendet mir wichtigen Trost?}
\vers{Dich ließ ich zurück, früh morgens bei Frost.}
\vers{Ich machte mich auf diese endlose Reise,}
\vers{doch ich komme zurück, ich wills dir beweisen!}
\end{strophe}

\phantom{a}

\begin{strophe}
\vers{Es sind nur so viele Streckenabschnitte...}
\vers{In mir stirbt ein Stück Hoffnung bei jedem der Schritte.}
\vers{Doch ich laufe, ich laufe, durch Dörfer und Felder,}
\vers{über Berge und Täler, zu dir immer ärger.}
\end{strophe}

\begin{strophe}
\vers{Ich kann die Strecke zwar schon erahnen --}
\vers{Oh nein, wie lange es dauert, bis die Dörfer sich weiten!}
\vers{Oh, wie lange es dauert bis die Berge sich scheiden!}
\end{strophe}

\begin{strophe}
\vers{Ich denke gerne zurück an des Sommers Gräser.}
\vers{Doch die Erinnerung verwelkt mir und wird immer blasser --}
\end{strophe}

\begin{strophe}
\vers{Ob du wohl an mich denkst, so wie ich an dich?}
\vers{Schwester, ich komme! Deine Haustür, brennt dort ein Licht?}
\vers{Ich will sie schon sehen! Sie ist so weit weg.}
\vers{Ach, beuge dich endlich, du elender Berg!}
\end{strophe}\\[0em]

%Sammlung der zehntausend Blätter ("Manyoshu") Teil 2, Gedicht 131 (万葉集－巻２－１３１) - Kakinomoto no Hitomaro (柿本人麻呂), ca. 653-710 n. Chr., kompiliert/im Manyoshu veröffentlicht um 759 n. Chr.
Sammlung der zehntausend Blätter ("Manyoshu") Teil 2,

Gedicht 131 - Kakinomoto no Hitomaro, ca. 653-710 n. Chr.,

kompiliert/im Manyoshu veröffentlicht um 759 n. Chr.\\[0em]

%Deutsche Übersetzung und lyrische Interpretation - Christoph Koloska mit Hilfe von verschiedenen Wörterbüchern (insb. 大辞林 第四版, 日本国語大辞典 精選版, 広辞苑 第七版), Google Übersetzer, DeepL und ChatGPT, 21.08.2026
Deutsche Übersetzung und lyrische Interpretation -

Christoph Koloska mit Hilfe von verschiedenen

Wörterbüchern (insb. Daijirin 4th ed., Nihonkokugodaijiten

carefully selected ed., Koujien 7th. ed.), Google

Übersetzer, DeepL und ChatGPT, 21.08.2026



